\section{机器人操作}

\subsection{灵巧操作的挑战}

灵巧操作（如用多指手抓取和操作物体）比移动更具挑战性：
\begin{itemize}
    \item 接触动力学高度非线性
    \item 状态空间维度更高
    \item 任务多样性更大
    \item Sim-to-real gap 在接触建模上更严重
\end{itemize}

\subsection{RL 在操作中的应用}

讲者展示了多个使用 RL 进行灵巧操作的案例，包括：
\begin{itemize}
    \item 旋转魔方（OpenAI Rubik's Cube）
    \item 笔旋转（pen spinning）
    \item 灵巧抓取（dexterous grasping）
\end{itemize}

这些任务都使用了大规模并行仿真 + domain randomization + teacher-student 的范式。

\subsection{本章小结}

RL 已经在多种机器人操作任务上展示了超越人类水平的能力，但仿真器质量仍是关键限制因素。

